\subsection{Compact and Cell Approximation}
\begin{theorem} \label{thm:compact-characterization-finite-outer-measure} \hlt{Compact Approximation at Finite Outer Measure}\\
Let $E\subseteq\R^p$ with $m_p^\ast(E)<+\infty$. Then $E\in\mathcal{L}^p$ iff $\forall\varepsilon>0$, $\exists$ compact $K\subseteq E$ with $m_p^\ast(E\setminus K)<\varepsilon$.
\begin{proof}
\begin{enumerate}[label=\roman*.]
\setlength{\itemsep}{0pt}
\item[$(\Rightarrow)$] Since $m(E)<+\infty$, choose closed $F\subseteq E$ with $m(E\setminus F)<\frac{\varepsilon}{2}$ by Theorem \ref{thm:closed-fsigma-characterization-lebesgue-measurable}.\\
Let $Q_N\defeq[-N,N]^p$. Since $F\cap Q_N\uparrow F$, continuity from below gives $N$ with $m(F\setminus Q_N)<\frac{\varepsilon}{2}$.\\
Then $K\defeq F\cap Q_N$ is compact, $K\subseteq E$, and $m_p^\ast(E\setminus K)=m(E\setminus K)\leq m(E\setminus F)+m(F\setminus Q_N)<\varepsilon$.
\item[$(\Leftarrow)$] Choose compact $K_n\subseteq E$ with $m_p^\ast(E\setminus K_n)<\frac{1}{n}$ and set $K\defeq\bigcup\limits_{n\geq1}K_n$.\\
Then $K$ is $F_\sigma$, hence $K\in\mathcal{L}^p$, and $E\setminus K\subseteq E\setminus K_n$ for all $n$, so $m_p^\ast(E\setminus K)=0$.\\
By Proposition \ref{prop:outer-null-sets-lebesgue-measurable}, $E\setminus K\in\mathcal{L}^p$, hence $E=K\cup(E\setminus K)\in\mathcal{L}^p$.
\end{enumerate}
\hlt{Structural Commentary}: Finite measure permits bounded truncation before inner closure.\\
\hlt{Key Insight}: Compact approximation is closed approximation plus a finite-measure tail estimate.
\end{proof}
\end{theorem}

\begin{theorem} \label{thm:finite-cell-symmetric-difference-approximation} \hlt{Finite Cell Approximation in Symmetric Difference}\\
If $E\in\mathcal{L}^p$ and $m(E)<+\infty$, then $\forall\varepsilon>0$, $\exists U=\bigcup\limits_{j=1}^{N}I_j$ with $I_j$ bounded open cells and $m(E\triangle U)<\varepsilon$.
\begin{proof}
Choose compact $K\subseteq E$ with $m(E\setminus K)<\frac{\varepsilon}{2}$ by Theorem \ref{thm:compact-characterization-finite-outer-measure}, and open $G\supseteq E$ with $m_p^\ast(G\setminus E)<\frac{\varepsilon}{2}$ by Theorem \ref{thm:open-gdelta-characterization-lebesgue-measurable}. For each $x\in K$, choose a bounded open cell $I_x$ with $x\in I_x$ and $\overline{I_x}\subseteq G$. Compactness gives $I_1,\ldots,I_N$ with $K\subseteq U\defeq\bigcup\limits_{j=1}^{N}I_j\subseteq G$. Then $E\triangle U\subseteq(E\setminus K)\cup(G\setminus E)$, so subadditivity gives $m(E\triangle U)<\varepsilon$.\\
\hlt{Structural Commentary}: Compactness turns a countable open cover into a finite geometric approximant.\\
\hlt{Key Insight}: Approximate from outside by cells and from inside by a compact core.
\end{proof}
\end{theorem}

\begin{corollary} \label{cor:cell-approximation-variant-choices} \hlt{Cell Approximation Variant Choices}\\
Let $E\in\mathcal{L}^p$ with $m(E)<+\infty$, and let $\varepsilon>0$. There exists $A\subseteq\R^p$ with $m(E\triangle A)<\varepsilon$ in each format:
\begin{enumerate}[label=(\roman*)]
\setlength{\itemsep}{0pt}
\item $A$ is a finite union of bounded open cells;
\item $A$ is a finite union of bounded closed cells;
\item $A$ is a finite union of bounded half-open cells of a fixed orientation.
\end{enumerate}
\begin{proof}
\begin{enumerate}[label=(\roman*)]
\setlength{\itemsep}{0pt}
\item By Theorem \ref{thm:finite-cell-symmetric-difference-approximation}, choose $A=\bigcup\limits_{j=1}^{N}I_j$ with $I_j$ bounded open cells and $m(E\triangle A)<\varepsilon$.
\item Choose a finite open-cell union $U=\bigcup\limits_{j=1}^{M}J_j$ with $m(E\triangle U)<\frac{\varepsilon}{2}$. If $M=0$, take $A=\emptyset$; otherwise choose closed cells $C_j\subseteq J_j$ with $m(J_j\setminus C_j)<\frac{\varepsilon}{2M}$ and set $A\defeq\bigcup\limits_{j=1}^{M}C_j$. Then $m(U\setminus A)<\frac{\varepsilon}{2}$, so $m(E\triangle A)<\varepsilon$.
\item Choose a finite open-cell union $U$ with $m(E\triangle U)<\varepsilon$. Refine by the finite coordinate grid from the endpoints of the open cells, and let $A$ be the union of the half-open grid atoms contained in $U$.\\
Then $U\triangle A$ is contained in finitely many cell faces, each outer-null, so $m(E\triangle A)<\varepsilon$.
\end{enumerate}
\hlt{Structural Commentary}: Cell format changes only boundary-null pieces.\\
\hlt{Key Insight}: A compact core inside an open envelope permits finite cell variants.
\end{proof}
\end{corollary}

\begin{corollary} \label{cor:real-line-finite-open-interval-approximation} \hlt{Real-Line Finite Open-Interval Approximation}\\
If $E\in\mathcal{L}^1$ and $m(E)<+\infty$, then $\forall\varepsilon>0$, $\exists O=\bigcup\limits_{j=1}^{N}I_j$ with $I_j$ pairwise disjoint bounded open intervals and $m(E\setminus O)+m(O\setminus E)=m(E\triangle O)<\varepsilon$.
\begin{proof}
By Theorem \ref{thm:finite-cell-symmetric-difference-approximation} with $p=1$, choose a finite union $U$ of bounded open intervals with $m(E\triangle U)<\varepsilon$. Write $U=O=\bigcup\limits_{j=1}^{N}I_j$ as the disjoint union of its open interval components.\\
Since $E\setminus O$ and $O\setminus E$ are disjoint measurable sets, additivity gives $m(E\triangle O)=m(E\setminus O)+m(O\setminus E)<\varepsilon$.\\
\hlt{Structural Commentary}: Merge finite open intervals into disjoint components.\\
\hlt{Key Insight}: A compact core makes the countable cover finite; interval merging makes it disjoint.
\end{proof}
\end{corollary}
